\section{Objetivos de Calidad}
Con el fin de asegurar un producto estable, seguro y confiable, se establecen los siguientes objetivos de calidad preventivos e integrales:

\begin{itemize}
    \item
  \textbf{Garantizar la Integridad y Consistencia de Datos:} El 100\% de
  los registros de incidencias deben asociarse a una ubicación
  normalizada en la estructura de base de datos (\texttt{pais},
  \texttt{provincia}, \texttt{ciudad}) y contar con coordenadas de
  latitud/longitud válidas, evitando redundancias u orfandad espacial.
\item
  \textbf{Minimizar Defectos Funcionales:} Lograr una tasa de aprobación
  de pruebas de al menos un 90\% en la validación de los flujos críticos
  (transición de estados, control de tiempos de resolución, asignación
  múltiple de roles y comentarios).
\item
  \textbf{Asegurar Rendimiento Temporal:} Garantizar que bajo
  condiciones normales de carga, los tiempos promedio de respuesta de la
  API REST del servidor y la carga inicial del mapa interactivo sean
  entre 3 y 5 segundos.
\item
  \textbf{Mitigar Riesgos de Seguridad Críticos:} Obtener una evaluación
  limpia de vulnerabilidades de seguridad web críticas, cumpliendo
  rigurosamente con los 5 vectores del estándar OWASP Top 10
  seleccionados.
\item
  \textbf{Mantener la Arquitectura y Código Limpio:} Implementar una
  arquitectura limpia y desacoplada mediante el uso de contenedores
  Docker (Nginx, Laravel, PostgreSQL), facilitando la mantenibilidad,
  extensibilidad y la futura integración de microservicios (por ejemplo,
  el módulo de IA/NLP o persistencia NoSQL).
\end{itemize}